\documentclass{article}
\usepackage{../../Self_Style}

\title{Phys 115A HW 6}
\author{Zih-Yu Hsieh}
\date{\today}

\begin{document}
\maketitle

\section*{1}

\begin{ques}\label{q1}
1 Hermitian Conjugates (Griffiths 3.5 + extra)

1. Find the Hermitian conjugates of x, i, and d/dx.

2. Show that for any two operators $\hat{Q}$ and $\hat{R}$ (not necessarily Hermitian), that $(\hat{Q}\hat{R})^\dagger = \hat{R}^\dagger \hat{Q}^\dagger$, $(\hat{Q}+ \hat{R})^\dagger = \hat{Q}^\dagger + \hat{R}^\dagger$, and $(c\hat{Q})^\dagger = c^\ast \hat{Q}^\dagger$ for a complex number $c$.

3. Construct the Hermitian conjugate of the harmonic oscillator raising operator $a^+$ (Eq. 2.48).

4. For a particle moving in one dimension, show that the operator $\hat{x}\hat{p}$ is not Hermitian. Can you construct an operator corresponding to some sort of physically observable product of position and momentum that is Hermitian?
\end{ques}

\begin{proof}

    \hfil

    \begin{enumerate}
        \item Here, if consider the Nuclear Space of integrable functions with domain $\RR$, codomain $\CC$, then given any such functions $\phi,\psi$, we have:
        \begin{align}
            \braket{\phi}{x\psi} &= \int_{-\infty}^{\infty}\phi^* \cdot (x\cdot \psi) dx = \int_{-\infty}^{\infty} (x\cdot \phi)^* \cdot \psi dx = \braket{x\phi}{\psi}
        \end{align}
        \begin{align}
            \braket{\phi}{i\psi} &= \int_{-\infty}^{\infty}\phi^* \cdot (i\psi)dx = \int_{-\infty}^{\infty}(-i\phi)^*\cdot \psi dx=\braket{-i\phi}{\psi}
        \end{align}
        \begin{align}
            \braket{\phi}{\frac{d}{dx}\psi} &= \int_{-\infty}^{\infty}\phi^* \cdot \frac{d}{dx}\psi dx = \phi^*\psi\bigg|_{-\infty}^{\infty}-\int_{-\infty}^{\infty}\left(\frac{d}{dx}\phi^*\right)\cdot \psi dx\\
            &= \int_{-\infty}^{\infty}\left(-\frac{d}{dx}\phi\right)^* \cdot \psi dx = \braket{-\frac{d}{dx}\phi}{\psi}
        \end{align}
        So, the above three relations show that $x^\dagger = x$, $i^\dagger = -i$, and $\left(\frac{d}{dx}\right)^\dagger = -\frac{d}{dx}$.

        \hfil

        \item Given any two operators $\hat{Q},\hat{R}$, let $\ket{\phi},\ket{\psi}$ be any two vectors, then:
        \begin{align}
            \braket{(\hat{Q}\hat{R})^\dagger \phi}{\psi}=\braket{\phi}{(\hat{Q}\hat{R})\psi} = \braket{\hat{Q}^\dagger \phi}{\hat{R}\psi} = \braket{\hat{R}^\dagger \hat{Q}^\dagger \phi}{\psi}
        \end{align}
        \begin{align}
            \braket{(\hat{Q}+\hat{R})^\dagger \phi}{\psi}&=\braket{\phi}{(\hat{Q}+\hat{R})\psi}=\braket{\phi}{\hat{Q}\psi}+\braket{\phi}{\hat{R}\psi}\\
            &= \braket{\hat{Q}^\dagger \phi}{\psi}+\braket{\hat{R}^\dagger \phi}{\psi} = \braket{(\hat{Q}^\dagger+\hat{R}^\dagger)\phi}{\psi}
        \end{align}
        \begin{align}
            \braket{(c\hat{Q})^\dagger \phi}{\psi}&= \braket{\phi}{c\hat{Q}\psi} = c\braket{\phi}{\hat{ Q}\psi} = c\braket{\hat{Q}^\dagger\phi}{\psi} = \braket{c^*\hat{Q}^\dagger \phi}{\psi}
        \end{align}
        Which, the above three relations show that $(\hat{Q}\hat{R})^\dagger = \hat{R}^\dagger \hat{Q}^\dagger$, $(\hat{Q}+\hat{R})^\dagger = \hat{Q}^\dagger+\hat{R}^\dagger$, and $(c\hat{Q})^\dagger = c^* \hat{Q}^\dagger$.

        \hfil

        \item Recall that the raising operator is defined as $a_+ = \frac{1}{\sqrt{2m\hbar w}}\left(x-i\frac{d}{dx}\right)$. So, using the result in part 1 and 2, we have:
        \begin{align}
            a_+^\dagger &= \left(\frac{1}{\sqrt{2m\hbar w}}\left(x-i\frac{d}{dx}\right)\right)^\dagger = \frac{1}{\sqrt{2m\hbar w}}\left(x^\dagger - \left(i\frac{d}{dx}\right)^\dagger\right) = \frac{1}{\sqrt{2m\hbar w}}\left(x + i\frac{d}{dx}\right) = a_-
        \end{align}
        Hence, the adjoint (Hermitian conjugate) of the raising operator is the lowering operator.

        \hfil

        \item Under one dimension, one has $\hat{x} \psi = x\psi$, and $\hat{p}\psi = -i\hbar\frac{d}{dx}\psi$. And, recall that both are self-adjoint/Hermitian operators (on the function space). So, we have the following:
        \begin{align}
            (\hat{x}\hat{p})^\dagger\psi &= \hat{p}\hat{x}\psi = -i\hbar\frac{d}{dx}(x\psi) = -i\hbar\psi + x\cdot \left(-i\hbar\frac{d}{dx}\psi\right) = \left(-i\hbar + \hat{x}\hat{p}\right)\psi
        \end{align}
        So, $(\hat{x}\hat{p})^\dagger = \hat{x}\hat{p} - i\hbar \neq \hat{x}\hat{p}$, showing it's not self-adjoint/Hermitian. However, if consider the operator $\frac{\hat{x}\hat{p}+\hat{p}\hat{x}}{2}$, since $(\hat{x}\hat{p})^\dagger = \hat{p}\hat{x}$ and $(\hat{p}\hat{x})^\dagger = \hat{x}\hat{p}$, we have:
        \begin{align}
            \left(\frac{\hat{x}\hat{p}+\hat{p}\hat{x}}{2}\right)^\dagger = \frac{(\hat{x}\hat{p})^\dagger+(\hat{p}\hat{x})^\dagger}{2}=\frac{\hat{p}\hat{x}+\hat{x}\hat{p}}{2}=\frac{\hat{x}\hat{p}+\hat{p}\hat{x}}{2}
        \end{align}
        Which is a desired self-adjoint/Hermitian operator corresponding to somewhat a product of position and momentum.
    \end{enumerate}
\end{proof}

\newpage

\section*{2}
\begin{ques}\label{q2}
2 Hilbert Space

1. Show that the set of all square-integrable functions (that is, the set of functions $f(x)$ for which $\int_a^b |f(x)|^2 dx$ is finite on a specified interval $x\in[a,b]$) is a vector space.

2. Show that $\langle f|g\rangle = \int_a^b f^\ast(x) g(x) dx$ satisfies the conditions for an inner product.
\end{ques}

\begin{proof}

    \hfil

    \begin{enumerate}
        \item Let $L$ denotes the set of all square-integrable functions on $[a,b]$. Then, for any $f,g \in L$ and any $z \in\CC$, we have:
        \begin{align}
            |f(x)+g(x)|^2 &\leq (|f(x)|+|g(x)|)^2 \leq (2\max\{|f(x)|,|g(x)|\})^2 = 4\max\{|f(x)|^2,|g(x)|^2\} \leq 4(|f(x)|^2+|g(x)|^2)
        \end{align}
        Hence, performing integration, one gets:
        \begin{align}
            \int_{a}^{b}|f(x)+g(x)|^2ds \leq 4\left(\int_{a}^{b}|f(x)|^2dx+\int_{a}^{b}|g(x)|^2dx\right) < \infty
        \end{align}
        Hence, $f(x)+g(x)\in L$, which $L$ is closed under addtion. Similarly, one has:
        \begin{align}
            \int_{a}^{b}|z f(x)|^2 dx = |z|^2\int_{a}^{b}|f(x)|^2 dx < \infty
        \end{align}
        Hence, $z f(x)\in L$, showing $L$ is closed under $\CC$-scalar multiplication.

        So, these two conditions conclude that $L$ is a $\CC$-vector space.

        \hfil

        \item To show the integral is an inner product, given any $f,g,h\in L$ and any $z, w\in\CC$, we have:
        \begin{align}
            \braket{f}{f}&=\int_{a}^{b}f(x)^*f(x)dx = \int_{a}^{b}|f(x)|^2dx \geq 0
        \end{align}
        \begin{align}
            \braket{f}{f}=0 \iff \int_{a}^{b}|f(x)|^2dx = 0\iff |f(x)|^2=0 \iff f(x)=0
        \end{align}
        (Note: for this one, it requires the continuity of $|f(x)|^2$ on $[a,b]$. For non-continuous functions, there ar pathological cases where $f(x)\neq 0$, yet such integral is $0$. For instance, $f(a)=1$ and $f(x)=0$ elsewhere works. In general, if the set of values where $f(x)\neq 0$ is measure $0$, then such integral is $0$).
        \begin{align}
            \braket{f}{z g+wh}=\int_{a}^{b}f(x)^*(zg(x)+w h(x))dx = z\int_{a}^{b}f(x)^*g(x)dx+w\int_{a}^{b}f(x)^*h(x)dx = z\braket{f}{g}+w\braket{f}{h}
        \end{align}
        \begin{align}
            \braket{g}{f}=\int_{a}^{b}g(x)^*f(x)dx = \int_{a}^{b}(f(x)^*g(x))^* dx = \left(\int_{a}^{b}f(x)^*g(x)dx\right)^* = \braket{f}{g}^*
        \end{align}
        So, these conditions characterize an inner product on $L$.
    \end{enumerate}
\end{proof}

\newpage

\section*{3 (part 3 a, c ND)}
\begin{ques}\label{q3}
3 Rigged Hilbert Space

1. Consider a Hilbert space $H$ that consists of all functions $\psi(x)$ such that $\int_{-\infty}^{\infty} |\psi(x)|^2 dx$ is finite. Show that there are functions in $H$ for which $\hat{x}\psi(x) \equiv x\psi(x)$ is not in $H$.

2. Consider the function space $\Omega$, called the nuclear space, which consists of all functions $\phi(x)$ that satisfy the infinite set of conditions $\int_{-\infty}^{\infty} |\phi(x)|^2 (1+|x|)^n dx < \infty$ for $n=0,1,2,\dots$. Show that for any $\phi(x)$ in $\Omega$, the function $\hat{x}\phi(x)\equiv x\phi(x)$ is also in $\Omega$.

3. Now consider the extended space $\Omega^\times$, which consists of all functions $\chi(x)$ that satisfy the condition $\langle \chi|\phi\rangle = \int_{-\infty}^{\infty} \chi^\ast(x)\phi(x)dx <\infty$ for all $\phi$ in the nuclear space $\Omega$. Which of the following functions belong to $\Omega$, to $H$, and/or to $\Omega^\times$?

(a) $\sin(x)$  
(b) $\sin(x)/x$  
(c) $x^2 \cos(x)$  
(d) $\dfrac{\log(1+|x|)}{1+|x|}$  
(e) $\exp(-x^2)$  
(f) $x^4 e^{-|x|}$
\end{ques}

\begin{proof}

    \hfil

    \begin{enumerate}
        \item Consider the function $\psi(x)=\frac{1}{\sqrt{x^2+1}}$, then $|\psi(x)|^2 = \frac{1}{x^2+1}$, showing $\int_{-\infty}^{\infty}|\psi(x)|^2dx = \arctan(x)\bigg|_{-\infty}^{\infty} = \pi$, hence $\psi(x)\in H$. However, if consider $x\psi(x)$, one has $|x\psi(x)|^2 = \frac{x^2}{x^2+1} = 1-\frac{1}{x^2+1}$, which the integral $\int_{a}^{b}\left(1-\frac{1}{x^2+1}\right)dx = (b-a) - \left(\arctan(x)\bigg|_{a}^{b}\right) = (b-a) - (\arctan(b)-\arctan(a))$. Which, take $b\rightarrow\infty$ and $a\rightarrow-\infty$, this integral clearly diverges (since $\arctan(b)-\arctan(a)$ is bounded, but not so for $b-a$). So, $x\psi(x)\notin H$.
        
        \hfil

        \item First, by induction one can show that $\int_{-\infty}^{\infty}|\phi(x)|^2\cdot|x|^n dx<\infty$ for all $n\in\NN$. For base case $n=0$, it is given by $\int_{-\infty}^{\infty}|\phi(x)|^2 \cdot |x|^0dx = \int_{-\infty}^{\infty}|\phi(x)|^2\cdot (1+|x|)^0 dx < \infty$. Now, suppose for all integer $k<n$ it satisfies $\int_{-\infty}^{\infty}|\phi(x)|^2 \cdot|x|^kdx < \infty$, then for $k=n$, one has the following:
        \begin{align}
            &|x|^n = (1+|x|)^n - \sum_{k=0}^{n-1}\begin{pmatrix}
                n\\k
            \end{pmatrix}|x|^k\\ 
            &\implies \int_{-\infty}^{\infty}|\phi(x)|^2 \cdot |x|^n dx = \int_{-\infty}^{\infty}|\phi(x)|^2 \cdot (1+|x|)^n dx - \sum_{k=0}^{n-1} \begin{pmatrix}
                n\\k
            \end{pmatrix}\int_{-\infty}^{\infty}|\phi(x)|^2 \cdot |x|^k dx < \infty
        \end{align}
        So, $\int_{-\infty}^{\infty}|\phi(x)|^2 \cdot |x|^n dx < \infty$, which completes the induction.

        So, this implies that given $\phi(x)\in\Sigma$, one has $x\phi(x)$ satisfies the following for all $n\in\NN$:
        \begin{align}
            \int_{-\infty}^{\infty}|x\phi(x)|^2 (1+|x|)^n dx = \sum_{k=0}^{n}\begin{pmatrix}
                n\\k
            \end{pmatrix}\int_{-\infty}^{\infty}|\phi(x)|^2 \cdot |x|^{k+2}dx < \infty
        \end{align}
        So, this indicates that $x\phi(x)\in \Omega^\times$ also.

        \hfil

        \item The spaces satisfy $\Omega \subseteq H \subseteq \Omega^\times$. Which, the functions given satisfy:
        \begin{itemize}
            \item[(a)] $\sin(x)\in\Omega^\times$: Given that $\int_{-\infty}^{\infty}|\sin(x)|^2 dx$ diverges (since integral of $\sin^2(x)$ diverges), then $\sin(x)\notin H$; however, since $-1\leq \sin(x)\leq 1$, any $\phi(x)\in \Omega^\times$ satisfies 
            
            \hfil

            \item[(b)] $\frac{\sin(x)}{x}\in H\setminus \Omega$, since one has $\int_{-\infty}^{\infty}\left|\frac{\sin(x)}{x}\right|^2 dx<\infty$ (since around $x=0$, one has $\int_{-1}^{1}\frac{\sin^2(x)}{x^2}dx$ converges due to the existence of limit at $x=0$, and the other two pieces $\int_{1}^{\infty}\frac{\sin^2(x)}{x^2}dx\leq \int_{1}^{\infty}\frac{1}{x^2}dx < \infty$, similar for the other piece on $(-\infty,-1)$).
            
            However, it's not in $\Omega$, since $x\cdot\frac{\sin(x)}{x}=\sin(x)\notin H$ (by part 3 (a)), hence $x\cdot\frac{\sin(x)}{x}\notin \Omega$, which it can't be in $\Omega$ (or else it contradicts the result in part 2).
            
            \hfil

            \item[(c)] d
            
            \hfil

            \item[(d)] $\frac{\log(1+|x|)}{1+|x|}\in H\setminus \Omega$. Recall that for all $\epsilon>0$, the growth order of $\log(z)$ is less than $z^\epsilon$. So, for simplicity one can pick $\epsilon=\frac{1}{4}$, so there exists $a\in\RR$, such that $\log(z) < z^\frac{1}{4}$ for all $a\in\RR$. As a result, one has the following:
            \begin{align}
                \int_{a}^{\infty}\left|\frac{\log(1+|x|)}{1+|x|}\right|^2dx < \int_{a}^{\infty}\frac{(1+|x|)^\frac{1}{2}}{(1+|x|)^2}dx = \int_{a}^{\infty}\frac{1}{(1+|x|)^\frac{3}{2}}dx < \infty
            \end{align}
            Which, $\int_{0}^{\infty}\left|\frac{\log(1+|x|)}{1+|x|}\right|^2dx<\infty$, with the function being symmetric, this implies $\int_{-\infty}^{\infty}\left|\frac{\log(1+|x|)}{1+|x|}\right|^2dx<\infty$.

            The reason it's not in $\Omega$, because take $n=2$, one has the following:
            \begin{align}
                \int_{-\infty}^{\infty}\left|\frac{\log(1+|x|)}{1+|x|}\right|^2 \cdot (1+|x|)^2 dx = \int_{-\infty}^{\infty}\log^2(1+|x|)dx \rightarrow \infty
            \end{align}
            Which, it doesn't satisfy the condition for being in $\Omega$.
            
            \hfil

            \item[(e)] $e^{-x^2}\in \Omega$, since for any $n\in\NN$, $(1+|x|)^n$ has growth order being a polynomial, so one has the following:
            \begin{align}
                \int_{-\infty}^{\infty}|e^{-x^2}|^2 (1+|x|)^ndx = \int_{-\infty}^{\infty}e^{-2x^2}(1+|x|)^n dx < \infty
            \end{align}
            Since $e^{-2x^2}$ has growth order much larger than the polynomials, the above integral converges.
            
            \hfil

            \item[(f)] $x^4e^{-|x|} \in \Omega$, since for any $n\in\NN$, one has $(1+|x|)^n$ being a polynomial, so is $|x|^8(1+|x|)^n$. As a result, one has:
            \begin{align}
                \int_{-\infty}^{\infty}|x^4e^{-|x|}|^2 (1+|x|)^n dx = \int_{-\infty}^{\infty}e^{-2|x|}\cdot |x|^8(1+|x|)^n dx < \infty
            \end{align}
            Since the improper integral on both sides are with exponential functions, which decay at a faster rate than any polynomials.
        \end{itemize}
    \end{enumerate}
\end{proof}

\newpage

\section*{4 (part 6 ND)}
\begin{ques}\label{q4}
4 Operators, Eigenvectors, Eigenvalues

What are the eigenvectors and eigenvalues for the following operators? In each case, how would you express a general initial wavefunction $\Psi(x,0)$ in terms of linear combinations of these eigenvectors?

1. The position operator $\hat{x}$ in position space.

2. The momentum operator $\hat{p}$ in position space.

3. The Hamiltonian for the free particle in position space.

4. The Hamiltonian for the infinite square well in position space.

5. The Hamiltonian for the harmonic oscillator in position space.

6. The lowering operator $a^-$ for the harmonic oscillator in position space.

7. The operator $i\frac{d}{d\phi}$, where $\phi$ is the polar angle in two dimensions, acting on the space of functions defined on the interval $0\le\phi\le 2\pi$ such that $f(\phi+2\pi)=f(\phi)$.

8. The operator $\frac{d^2}{d\phi^2}$, acting on the same space.
\end{ques}

\begin{proof}

    \hfil

    \begin{enumerate}
        \item Given $\hat{x}=x$, the eigenvectors are $\delta(x-a)$, corresponding to eigenvalue $a\in\RR$ (since $x\delta(x-a) = a\delta(x-a)$ as operators when integrating with general functions). For general initial wavefunction $\Psi(x,0)$, one can express it as:
        \begin{align}
            \Psi(x,0) = \int_{-\infty}^{\infty}\Psi(a,0)\delta(x-a)da
        \end{align}

        \hfil

        \item Given $\hat{p} = -i\hbar\frac{d}{dx}$, the eigenvectors and eigenvalues are $\frac{1}{\sqrt{2\pi}}e^{ikx}$, with eigenvalue $\hbar k$ for $k\in\RR$ (since taking derivative one yields $\frac{ik}{\sqrt{2\pi}}e^{ikx}$, multiply $-i\hbar$ gives $\frac{\hbar k}{\sqrt{2\pi}}e^{ikx}$). For general initial wavefunction $\Psi(x,0)$, one can express it as (using Fourier / Inverse Fourier Transform):
        \begin{align}
            \Psi(x,0) = \frac{1}{\sqrt{2\pi}}\int_{-\infty}^{\infty}\phi(k)e^{ikx}dk,\quad \phi(k) = \frac{1}{\sqrt{2\pi}}\int_{-\infty}^{\infty}\Psi(x,0)e^{-ikx}dx
        \end{align}

        \hfil

        \item Given Hamiltonia for the free partical $\hat{H}=-\frac{\hbar^2}{2m}\frac{d^2}{dx^2}$, the eigenvectors are again $\frac{1}{\sqrt{2\pi}}e^{ikx}$ for $k\in\RR$, with eigenvalue $\frac{k^2\hbar^2}{2m}$ (since after taking second derivative, one yields $-\frac{k^2}{\sqrt{2\pi}}e^{ikx}$, multiply by $-\frac{\hbar^2}{2m}$ gives $\frac{k^2\hbar^2}{2m}\cdot\frac{1}{\sqrt{2\pi}}e^{ikx}$). For general initial wavefunction $\Psi(x,0)$, one again express it as follow, similar to part 2:
        \begin{align}
            \Psi(x,0) = \frac{1}{\sqrt{2\pi}}\int_{-\infty}^{\infty}\phi(k)e^{ikx}dk,\quad \phi(k) = \frac{1}{\sqrt{2\pi}}\int_{-\infty}^{\infty}\Psi(x,0)e^{-ikx}dx
        \end{align}

        \hfil

        \item The Hamiltonian for the infinite square well is $\hat{H}=-\frac{\hbar^2}{2m}\frac{d^2}{dx^2}$, where $x\in [0,a]$, with condition $\psi(0)=\psi(a)=0$, and $\psi(x)=0$ for $x\notin [0,a]$. The eigenvectors are $\sqrt{\frac{2}{a}}\sin\left(\frac{n\pi}{a}x\right)$ ($x\in[0,a]$) for $n\in\ZZ^{\geq 0}$, with eigenvalue $E_n = \frac{1}{2m}\left(\frac{n\pi\hbar}{a}\right)^2$ (since after taking second derivative, one yields $-\frac{n^2\pi^2}{a^2}\cdot\sqrt{\frac{2}{a}}\sin\left(\frac{n\pi}{a}x\right)$, multiply by $-\frac{\hbar^2}{2m}$ gives $\frac{n^2\pi^2\hbar^2}{2ma^2}\sqrt{\frac{2}{a}}\sin\left(\frac{n\pi}{a}x\right)$). For general initial wavefunction $\Psi(x,0)$ can be expressed as follow:
        \begin{align}
            \Psi(x,0)=\sum_{n=1}^{\infty}c_n \sqrt{\frac{2}{a}}\sin\left(\frac{n\pi}{a}x\right),\quad c_n = \int_{0}^{a}\Psi(x,0)\sqrt{\frac{2}{a}}\sin\left(\frac{n\pi}{a}x\right)dx
        \end{align}

        \hfil

        \item The Hamiltonian for harmonic oscillator is $\hat{H}=-\frac{\hbar^2}{2m}\frac{d^2}{dx^2}+\frac{1}{2}mw^2x^2$, so the eigenvectors are given by $\psi_n(x) = \frac{1}{\sqrt{n!}}(a_+)^n\psi_0(x)$, where $\psi_0(x)=\left(\frac{mw}{\pi\hbar}\right)^\frac{1}{4}e^{-\frac{mw}{2\hbar}x^2}$ and the operator $a_+ = \frac{1}{\sqrt{2mw\hbar}}(\hat{x}-i\hat{p})$, with eigenvalue $E_n = \hbar w\left(n+\frac{1}{2}\right)$ (these are the results from the Harmonic Oscillator chapter). For general initial wavefunction $\Psi(x,0)$, one can express it as:
        \begin{align}
            \Psi(x,0) = \sum_{n=0}^{\infty}c_n \psi_n(x),\quad c_n=\int_{-\infty}^{\infty}\Psi(x,0)\psi_n(x)dx
        \end{align}

        \hfil

        \item Given the lowering operator $a_- = \frac{1}{\sqrt{2mw\hbar}}(mw\hat{x}+i\hat{p}) = \frac{1}{\sqrt{2mw\hbar}}\left(mwx + \hbar\frac{d}{dx}\right)$. Which, suppose $\lambda\in\CC$ and $\psi:\RR\rightarrow\CC$ satisfies the following:
        \begin{align}
            a_- \psi = \lambda \psi &\implies mwx\psi + \hbar\frac{d\psi}{dx} = \sqrt{2mw\hbar} \lambda \psi\\
            &\implies \frac{d\psi}{dx} +\left(\frac{mw}{\hbar}x-\sqrt{\frac{2mw}{\hbar}}\lambda\right)\psi=0
        \end{align}
        Which, using integrating factor, we get the following function as eigenvector, with eigenvalue $\lambda$:
        \begin{align}
            &I = \exp\left(\int \left(\frac{mw}{\hbar}x - \sqrt{\frac{2mw}{\hbar}}\lambda\right)dx\right) = \exp\left(\frac{mw}{2\hbar}x^2 - \sqrt{\frac{2mw}{\hbar}}\lambda x\right)\\
            &\implies I \psi = C\implies \psi = CI^{-1} = C\exp\left(-\frac{mw}{2\hbar}x^2 + \sqrt{\frac{2mw}{\hbar}}\lambda x\right)
        \end{align}
        Which, take $ik := \sqrt{\frac{2mw}{\hbar}}\lambda$, any $\phi(k)$ that converges under Fourier Transform satisfies:
        \begin{align}
            
        \end{align}
        
        \hfil

        \item Given $f(\phi+2\pi)=f(\phi)$, if it's the eigenvector of $i\frac{d}{d\phi}$ with eigenvalue $\lambda$, one has $i\frac{d}{d\phi}f = \lambda f$, or $\frac{df}{d\phi}=-i\lambda f$, showing $f(\phi) = Ke^{-i\lambda \phi}$. As a result, with period being multiples of $2\pi$, it enforces $\lambda = n$ for $n\in\ZZ$. So, the eigenvectors can be chosen as $\frac{1}{\sqrt{2\pi}}e^{in\phi}$ for $n\in\ZZ$, with eigenvalue $-n$. For general initial wavefunction $\Psi(\phi,0)$, one can express it as follow (using Fourier series):
        \begin{align}
            \Psi(x,0) = \sum_{n\in\ZZ}\frac{c_n}{\sqrt{2\pi}}e^{in\phi}, \quad c_n = \int_{0}^{2\pi}\Psi(\phi,0)\frac{1}{\sqrt{2\pi}}e^{-in\phi}d\phi
        \end{align}

        \hfil

        \item Given $f$ as an eigenvector of $\frac{d^2}{d\phi^2}$ with eigenvalue $-\lambda^2$, one has $\frac{d^2f}{d\phi^2}=-\lambda^2\phi$, showing $f(\phi) = Ke^{i\lambda\phi} + Le^{-i\lambda\phi}$. Again, with period being multiples of $2\pi$, this again enforces $\lambda = n$ for $n\in\ZZ$. So, the eigenvectors can be collected to $\frac{1}{\sqrt{2\pi}}e^{in\phi}$ again, with eigenvalue $-n^2$. For general inital wavefunction $\Psi(\phi,0)$, one can again express it identically as part 7:
        \begin{align}
            \Psi(x,0) = \sum_{n\in\ZZ}\frac{c_n}{\sqrt{2\pi}}e^{in\phi}, \quad c_n = \int_{0}^{2\pi}\Psi(\phi,0)\frac{1}{\sqrt{2\pi}}e^{-in\phi}d\phi
        \end{align}
    \end{enumerate}
\end{proof}

\newpage

\section*{5}
\begin{ques}\label{q5}
5 (Your Name) Uncertainty Principle (Griffiths 3.15)

Using the generalized uncertainty principle, prove the “(Your Name) Uncertainty Principle’’ relating the uncertainty in position $\sigma_x$ to the uncertainty in energy $\sigma_H$:
\[
\sigma_x \sigma_H \ge \frac{\hbar}{2m} |\langle p\rangle|
\]
for $H = \hat{p}^2/2m + V(\hat{x})$. For stationary states this is not particularly helpful – why not?
\end{ques}

\begin{proof}
    Given any two operators $\hat{A},\hat{B}$ corresponding to observables $A,B$, the generalized uncertainty principle states:
    \begin{align}
        \sigma_A\sigma_B \geq \frac{1}{2}\left|\left<[\hat{A},\hat{B}]\right>\right|
    \end{align}
    Here, choose $A=x$ and $B=H$, then $[\hat{x},\hat{H}]$ can be given as follow:
    \begin{align}
        [\hat{x},\hat{H}]f &= x\left(-\frac{\hbar^2}{2m}\frac{d^2}{dx^2}+V(x)\right)f - \left(-\frac{\hbar^2}{2m}\frac{d^2}{dx^2}+V(x)\right)(xf)\\
        &= -x\cdot \frac{\hbar^2}{2m}\frac{d^2f}{dx^2} + xV(x)f + \frac{\hbar^2}{2m}\frac{d^2}{dx^2}(xf) - xV(x)f\\
        &= -\frac{\hbar^2}{2m}\cdot x\frac{d^2f}{dx^2} + \frac{\hbar^2}{2m}\left(2\frac{df}{dx}+x\frac{d^2f}{dx^2}\right)\\
        &= \frac{i\hbar}{m}\left(-i\hbar\frac{d}{dx}\right)f = -\frac{i\hbar}{m}\hat{p}f
    \end{align}
    So, it's concluded that on the function space, $[\hat{x},\hat{H}] = -\frac{i\hbar}{m}\hat{p}$. Hence, applying the generalized uncertainty principle, one has:
    \begin{align}
        \sigma_x\sigma_H \geq \frac{1}{2}\left|\left<[\hat{x},\hat{H}]\right>\right| = \frac{1}{2}\left|\left<-\frac{i\hbar}{m}\hat{p}\right>\right| = \frac{\hbar}{2m}\left|\left<\hat{p}\right>\right|
    \end{align}

    For stationary state, this is not usefule since $\sigma_H = 0$ (as it is an eigenvector of $\hat{H}$), so the inequality is enforces to be $0\geq \frac{\hbar}{2m}\left|\left<p\right>\right|$, or enforcing $\left<p\right>=0$. This doesn't provide further information about the state's uncertainty in position.
\end{proof} 

\newpage

\section*{6}
\begin{ques}\label{q6}
6 Spectra of Operators (Griffiths 3.16 + extra)

1. Show that two noncommuting operators cannot have a complete set of common eigenfunctions.

2. Consider two operators $\hat{P},\hat{Q}$ for which $[\hat{P},\hat{Q}]=0$. $\hat{P}$ has a spectrum of non-degenerate eigenvectors $\psi_i$ with eigenvalues $p_i$. You prepare a particle. A measurement of $P$ yields value $p_5$. If you measure $Q$ afterwards, what do you find?

3. $\hat{Q}$ has only non-degenerate eigenvectors $\psi_i$ and eigenvalues $Q_i$. What are the eigenvectors and eigenvalues of the inverse operator $\hat{Q}^{-1}$?
\end{ques}

\begin{proof}
    
    \hfil

    \begin{enumerate}
        \item We'll prove the contrapositive: Suppose $\hat{P},\hat{Q}$ are two operators with complete set of common eigenvectors $\{\psi_i\}_{i\in I}$, each associated with eigenvalue $\{p_i\}_{i\in I}$ for $\hat{P}$ and $\{q_i\}_{i\in I}$ for $\hat{Q}$ (or, the eigenvectors span the whole space). Then, for any function $f$, one can express it as $f=\sum_{i \in I}a_i\psi_i$, where $a_i\in\CC$ are chosen so that this converges. So, consider the action of $[\hat{P},\hat{Q}]$, we have:
        \begin{align}
            [\hat{P},\hat{Q}]f = \sum_{i\in I}a_i(\hat{P}\hat{Q}-\hat{Q}\hat{P})\psi_i = \sum_{i\in I}a_i(p_i q_i-q_i p_i)\psi_i=0
        \end{align}
        (Note: since $\hat{Q}\psi_i = q_i\psi_i$, and $\hat{P}=p_i\psi_i$).

        Hence, on the desired vector space, $[\hat{P},\hat{Q}]=0$, showing $\hat{P},\hat{Q}$ commutes.

        \hfil

        \item Since $[\hat{P},\hat{Q}]=0$, one has $\hat{P}\hat{Q}=\hat{Q}\hat{P}$. Which, applying on the eigenvectors $\psi_i$, one has:
        \begin{align}
            \hat{P}(\hat{Q}\psi_i) = \hat{Q}(\hat{P}\psi_i) = p_i \hat{Q}\psi_i
        \end{align}
        Which, let $E(\hat{P},p_i)$ denotes the eigenspace of $\hat{P}$ with eigenvalue $p_1$, then $\hat{Q}\psi_i \in E(\hat{P},p_i)$ (showing $\hat{Q}$ maps the eigenspace to itself). 

        Hence, if measurement of $P$ yields value $p_5$, the state is currently at $\psi_5$; which, after measuring $Q$, one has $\hat{Q}\psi_t\in E(\hat{P},p_5)$, showing that if measuring $P$ again, one definitively gets $p_5$. (Even though the state may change, but it's still within the sam eigenspace of $\hat{P}$, so the measured observable of $P$ stays the same).

        \hfil

        \item The eigenvectors and eigenvalues of $\hat{Q}^{-1}$ are precisely given by $\psi_i$ and $\frac{1}{Q_i}$ respectively. 
        
        First, the $\psi_i$, $\frac{1}{Q_i}$ are eigenvectors / eigenvalues of $\hat{Q}^{-1}$, since:
        \begin{align}
            \psi_i=\hat{Q}^{-1}(\hat{Q}\psi_i) = \hat{Q}^{-1}(Q_i\psi_i)\implies \hat{Q}^{-1}\psi_i = \frac{1}{Q_i}\psi_i
        \end{align}
        Conversely, if $v, \frac{1}{\lambda}$ are eigenvector / eigenvalue of $\hat{Q}^{-1}$, then:
        \begin{align}
            v = \hat{Q}(\hat{Q}^{-1}v) = \hat{Q}\left(\frac{1}{\lambda} v\right)\implies \hat{Q}v = \lambda v
        \end{align}S
        o, $v,\lambda$ are eigenvector / eigenvalue of $\hat{Q}$, by the claim one must have $v=\psi_i$ and $\lambda=Q_i$ for some $i$. So, the eigenvector / eigenvalue of $\hat{Q}^{-1}$ must be $\psi_i$ and $\frac{1}{Q_i}$ respectively.

    \end{enumerate}
\end{proof}

\newpage

\section*{7}
\begin{ques}\label{q7}
7 Sequential Measurements (Griffiths 3.33)

Consider an observable $A$ with eigenvectors $\psi_1,\psi_2$ and eigenvalues $a_1,a_2$. Consider a second observable $B$ with eigenvectors $\phi_1,\phi_2$ and eigenvalues $b_1,b_2$. The TA tells you:
\[
\psi_1 = \frac{1}{5}(3\phi_1 + 4\phi_2), \qquad
\psi_2 = \frac{1}{5}(4\phi_1 - 3\phi_2)
\]

1. You prepare a particle and measure $A$, getting $a_1$. What is the state afterwards?

2. Then measure $B$. What are the results and their probabilities?

3. Immediately measure $A$ again. What is the probability of getting $a_1$?
\end{ques}

\begin{proof}

    \hfil

    \begin{enumerate}
        \item If measure observable $A$ and get $a_1$, it implies the state collapses to the eigenvector of $A$ corresponding to $a_1$, which is $\psi_1$.
        
        \hfil

        \item Since after measuring $A$, the current state is $\psi_1=\frac{1}{5}(3\phi_1+4\phi_2)$ (in terms of the eigenvectors of $B$), so measuring $B$, the possible results are the eigenvalues of $B$, which are $b_1,b_2$. 
        
        And, since in this state $\phi_1$ (the eigenvector of $b_1$) has coefficient $\frac{3}{5}$, while $\phi_2$ (the eigenvector of $b_2$) has coefficient $\frac{4}{5}$, the probabilities of getting $b_1,b_2$ are $\PP(b_1)=\left|\frac{3}{5}\right|^2 = \frac{9}{25} = 36\%$, and $\PP(b_2)=\left|\frac{4}{5}\right|^2 = \frac{16}{25}=64\%$.

        \hfil

        \item First, we'll express $\phi_1,\phi_2$ (the eigenvectors of $B$, corresponding to eigenvalues $b_1,b_2$) in terms of $\psi_1,\psi_2$ (the eigenvectors of $A$). Which, based on the given relation, by scalar multiplication we get:
        \begin{align}
            &20\psi_1 = 12\phi_1 + 16\phi_2,\quad 15\psi_2 = 12\phi_1 - 9\phi_2\\
            &\implies 20\psi_1 - 15\psi_2 = 25\phi_2,\quad \phi_2 = \frac{1}{5}(4\psi_1 - 3\psi_2)\\
        \end{align}
        Again, with scalar multiplication and substitution, we have:
        \begin{align}
            &25\psi_1 = 15\phi_1 + 20\phi_2 = 15\phi_1+16\psi_1 - 12\psi_2\\
            &\implies 15\phi_1 = 9\psi_1 + 12\psi_2,\quad \phi_1 = \frac{1}{5}(3\psi_1+4\psi_2)
        \end{align}
        Since it's possible that the state becomes $\phi_1$ (with probability $36\%$) or $\phi_2$ (with probability $64\%$). After measuring $A$, the probability of getting $a_1$ is given by the conditional probability. 

        If ending with $b_1$ previously (or at state $\phi_1$), since $\psi_1$ (the eigenvector of $a_1$) has $\phi_1$ coefficient $\frac{3}{5}$, the probability of getting $a_1$ is $\PP(a_1|b_1)=\left|\frac{3}{5}\right|^2 = \frac{9}{25}=36\%$; else, if ending with $b_2$ previously (or at state $\phi_2$), since $\psi_1$ has $\phi_1$ coefficient $\frac{4}{5}$, the probability of getting $a_1$ is $\PP(a_1|b_2)\left|\frac{4}{5}\right|^2 = 64\%$. So, the total probability of measuring $a_1$ is:
        \begin{align}
            \PP(a_1) &= \PP(a_1|b_1)\cdot\PP(b_1) + \PP(a_1|b_2)\cdot\PP(b_2) = \left(\frac{9}{25}\right)^2+\left(\frac{16}{25}\right)^2=\frac{337}{625} \approx 53.92\%
        \end{align}
    \end{enumerate}
\end{proof}

\newpage

\section*{8}
\begin{ques}\label{q8}
8 Harmonic Oscillator in “energy space’’ (Griffiths 3.39)

1. Compute $\langle n|\hat{x}|n'\rangle$ for general $n,n'$.

2. Compute $\langle n|\hat{p}|n'\rangle$.

3. Construct the infinite-dimensional matrix representation of $\hat{x}$ using these matrix elements.

4. Construct the infinite-dimensional matrix for $\hat{p}$.

5. Show that $H = \frac{1}{2m}p^2 + \frac{1}{2}m\omega^2 x^2$ is diagonal in this basis. What are the diagonal elements?
\end{ques}

\begin{proof}

    First, recall that in terms of ladder operators, $\hat{x}=\sqrt{\frac{\hbar}{2mw}}(a_++a_-)$, and $ \hat{p}=i\sqrt{\frac{\hbar mw}{2}}(a_+-a_-)$. 
    
    Also, if $\ket{n'}$ represents the eigenvector of the harmonic oscillator Hamiltonian, then $a_+\ket{n'}=\sqrt{n'+1}\ket{n'+1}$, and $a_-\ket{n'}=\sqrt{n'}\ket{n'-1}$ (recall the relation $a_+\psi_n(x)=\sqrt{n+1}\psi_{n+1}(x)$, and $a_-\psi_n(x)=\sqrt{n}\psi_{n-1}(x)$). 

    Finally, we also have $a_-\ket{0} = 0$ (due to the fact $a_-\psi_0(x)=0$).

    \begin{enumerate}
        \item Given the above relations, we have the following for all $n\in \NN$, and $n'\in\ZZ_{>0}$:
        \begin{align}
            \hat{x}\ket{n'} = \sqrt{\frac{\hbar}{2mw}}\left(\sqrt{n'+1}\ket{n'+1}+\sqrt{n
            }\ket{n'-1}\right)\implies \bra{n}\hat{x}\ket{n'} = \begin{cases}
                \sqrt{\frac{\hbar(n'+1)}{2mw}} & n=n'+1\\
                \sqrt{\frac{\hbar n'}{2mw}} & n=n'-1\\
                0 & \text{otherwise}
            \end{cases}
        \end{align}
        Which, a special case happens when $n'=0$, where $\hat{x}\ket{0} =\sqrt{\frac{\hbar}{2mw}}\ket{1}$, so $\bra{n}\hat{x}\ket{0} = \begin{cases}
            \sqrt{\frac{\hbar}{2mw}} & n=1\\
            0 & \text{otherwise}
        \end{cases}$.

        \hfil

        \item We also have the following for all $n\in\NN$ and $n'\in\ZZ_{>0}$:
        \begin{align}
            \hat{p}\ket{n'} = i\sqrt{\frac{\hbar mw}{2}}(\sqrt{n'+1}\ket{n'+1}-\sqrt{n'}\ket{n'-1})\implies \bra{n}\hat{p}\ket{n'} = \begin{cases}
                i\sqrt{\frac{\hbar mw(n'+1)}{2}} & n=n'+1\\
                -i\sqrt{\frac{\hbar mw n'}{2}} & n=n'-1\\
                0 & \text{otherwise}
            \end{cases}
        \end{align}
        Again, a special case happens when $n'=0$, where $\hat{p}\ket{0} = i\sqrt{\frac{\hbar mw}{2}}\ket{1}$, so $\bra{n}\hat{p}\ket{0} = \begin{cases}
            i\sqrt{\frac{\hbar mw}{2}} & n=1\\
            0 & \text{otherwise}
        \end{cases}$.

        \hfil

        \item Based on the results in part 1, if $\mathcal{M}(\hat{x}) = (a_{ij})_{0\leq i,j}$ represents the matrix under this basis, then one has:
        \begin{align}
            a_{ij}=\begin{cases}
                0 & i=j \text{ or } |i-j|\geq 2\\
                \sqrt{\frac{\hbar (j+1)}{2mw}} & i=j+1\\
                \sqrt{\frac{\hbar j}{2mw}} & i=j-1
            \end{cases}
        \end{align}
        Writing out the first several rows and columns, we have:
        \begin{align}
            \sqrt{\frac{\hbar}{2mw}}\begin{pmatrix}
                0 & 1&0&0&0&\cdots\\
                1 & 0&\sqrt{2}&0&0&\cdots\\
                0&\sqrt{2}&0&\sqrt{3}&0&\cdots\\
                0 & 0& \sqrt{3} & 0 & \sqrt{4} & \cdots\\
                0 & 0 & 0 & \sqrt{4} & 0 & \cdots\\
                \vdots & \vdots & \vdots & \vdots & \vdots & \ddots
            \end{pmatrix}
        \end{align}

        \hfil

        \item Based on th results in part 2, if $\mathcal{M}(\hat{p}) = (b_{jk})_{0\leq j,k}$ represents the matrix under this basis, then one has:
        \begin{align}
            b_{jk} = \begin{cases}
                0 & j=k \text{ or } |j-k|\geq 2\\
                i\sqrt{\frac{\hbar mw(k+1)}{2}} & j=k+1\\
                -i\sqrt{\frac{\hbar mwk}{2}} & j=k-1
            \end{cases}
        \end{align}
        Writing out the first several rows and columns, we have:
        \begin{align}
            i\sqrt{\frac{\hbar mw}{2}}\begin{pmatrix}
                0 & -1 & 0&0&0&\cdots\\
                1&0&-\sqrt{2}&0&0&\cdots\\
                0&\sqrt{2}&0&-\sqrt{3}&0&\cdots\\
                0&0&\sqrt{3}&0&-\sqrt{4}&\cdots\\
                0&0&0&\sqrt{4}&0&\cdots\\
                \vdots &\vdots &\vdots &\vdots &\vdots &\ddots
            \end{pmatrix}
        \end{align}

        \hfil

        \item To show $\hat{H}=\frac{\hat{p}^2}{2m}+\frac{1}{2}mw^2\hat{x}^2$ is diagonal in this basis, we first yield the following computation for all $n'\geq 2$:
        \begin{align}
            \hat{p}^2\ket{n'} &= \hat{p} i\sqrt{\frac{\hbar mw}{2}}(\sqrt{n'+1}\ket{n'+1} - \sqrt{n'}\ket{n'-1})\\
            &= -\frac{\hbar mw}{2}(\sqrt{(n'+1)(n'+2)}\ket{n'+2} - (n'+1)\ket{n'} - n'\ket{n'}+ \sqrt{n'(n'-1)}\ket{n'-2})\\
            &= -\frac{\hbar mw}{2}(\sqrt{(n'+1)(n'+2)}\ket{n'+2} + \sqrt{n'(n'-1)}\ket{n'-2} - (2n'+1)\ket{n'})
        \end{align}
        \begin{align}
            \hat{x}^2 \ket{n'} &= \hat{x}\sqrt{\frac{\hbar}{2mw}}(\sqrt{n'+1}\ket{n'+1} + \sqrt{n'}\ket{n'-1})\\
            &= \frac{\hbar}{2mw}(\sqrt{(n'+1)(n'+2)}\ket{n'+2} + (n'+1)\ket{n'} + n'\ket{n'} + \sqrt{n'(n'-1)}\ket{n'-2})\\
            &= \frac{\hbar}{2mw}(\sqrt{(n'+1)(n'+2)}\ket{n'+2} + (2n'+1)\ket{n'} + \sqrt{n'(n'-1)}\ket{n'-2})
        \end{align}
        Which, the action of $\hat{H}$ is as follow:
        \begin{align}
            \hat{H}\ket{n'} =& \frac{1}{2m}\hat{p}^2\ket{n'}+\frac{1}{2}mw^2\hat{x}^2\ket{n'}\\
            =& -\frac{\hbar w}{4}(\sqrt{(n'+1)(n'+2)}\ket{n'+2}+\sqrt{n'(n'-1)}\ket{n'-2}-(2n'+1)\ket{n'})\\
            &+ \frac{\hbar w}{4}(\sqrt{(n'+1)(n'+2)}\ket{n'+2}+\sqrt{n'(n'-1)}\ket{n'-2}+(2n'+1)\ket{n'})\\
            =& \frac{\hbar w}{2}(2n'+1)\ket{n'} = \hbar w\left(n'+\frac{1}{2}\right)\ket{n'}
        \end{align}

        \hfil

        For special case $n=0$, one has:
        \begin{align}
            &\hat{p}^2\ket{0} = \hat{p}i\sqrt{\frac{\hbar mw}{2}}\ket{1} = -\frac{\hbar mw}{2}(\sqrt{2}\ket{2}-\ket{0})\\
            &\hat{x}^2\ket{0} = \hat{x}\sqrt{\frac{\hbar}{2mw}}\ket{1} = \frac{\hbar}{2mw}(\sqrt{2}\ket{2}+\ket{0})\\
            &\implies \hat{H}\ket{0} = \frac{1}{2m}\hat{p}^2\ket{0}+\frac{1}{2}mw^2\hat{x}^2\ket{0} = -\frac{\hbar w}{4}(\sqrt{2}\ket{2}-\ket{0}) + \frac{\hbar w}{4}(\sqrt{2}\ket{2}+\ket{0}) = \frac{\hbar w}{2}\ket{0}
        \end{align}

        Also, for special case $n=1$, one has:
        \begin{align}
            &\hat{p}^2\ket{1} = \hat{p}i\sqrt{\frac{\hbar mw}{2}}(\sqrt{2}\ket{2}-\ket{0})= -\frac{\hbar mw}{2}(\sqrt{6}\ket{3} - 3\ket{1})\\
            &\hat{x}^2\ket{1} = \hat{x}\sqrt{\frac{\hbar}{2mw}}(\sqrt{2}\ket{2}+\ket{0}) = \frac{\hbar}{2mw}(\sqrt{6}\ket{3}+3\ket{1})\\
            &\implies \hat{H}\ket{1} = \frac{1}{2m}\hat{p}^2\ket{1}+\frac{1}{2}mw^2\hat{x}^2\ket{1} = -\frac{\hbar w}{4}(\sqrt{6}\ket{3}-3\ket{1}) + \frac{\hbar w}{4}(\sqrt{6}\ket{3}+3\ket{1}) = \frac{3\hbar w}{2}\ket{1}
        \end{align}
        Which, one can see the action of $\hat{H}$ on $\ket{n'}$ is by scaling with factor $E_n = \hbar w\left(n'+\frac{1}{2}\right)$, which is exactly the eigenvalue / energy of the harmonic oscillator statopmaru states.
    \end{enumerate}
\end{proof}

\newpage

\section*{9}
\begin{ques}\label{q9}
9 A two-level system

The Hamiltonian in basis $|1\rangle,|2\rangle$ is
\[
\hat{H} = \varepsilon(-|1\rangle\langle1| + |2\rangle\langle2| + |1\rangle\langle2| + |2\rangle\langle1|)
\]

1. Find the eigenvalues and eigenvectors.

2. Write the matrix representation of $\hat{H}$ in this basis.
\end{ques}

\begin{proof}

    If input $\ket{1}, \ket{2}$ respectively, we have $\hat{H}\ket{1} = -\epsilon\ket{1} + \epsilon\ket{2}$, and $\hat{H}\ket{2} = \epsilon \ket{1}+\epsilon\ket{2}$ (based on the orthonormality $\braket{1}=\braket{2}=1$, and $\braket{1}{2}=\braket{2}{1}=0$). Here, we'll assume $\epsilon\neq 0$.

    \begin{enumerate}
        \item Suppose $\lambda\in\CC$ is an eigenvalue of $\hat{H}$ and vector $a\ket{1}+b\ket{2}$ is an eigenvector corresponding to it. Then, we have:
        \begin{align}
            \lambda(a\ket{1}+b\ket{2}) &=\hat{H}(a\ket{1}+b\ket{2})=a(-\epsilon\ket{1}+\epsilon\ket{2})+b(\epsilon\ket{1}+\epsilon\ket{2}) = \epsilon(b-a)\ket{1}+\epsilon(b+a)\ket{2}
        \end{align}
        So, $\lambda a=\epsilon(b-a)$, and $\lambda b=\epsilon(b+a)$.

        First, suppose $a\neq 0$, WLOG can assume $a=1$ up to scaling, then $\lambda=\epsilon(b-1)$, plugin to the second equation provides $\epsilon(b-1)b=\epsilon(b+1)$, or $b^2-2b-1=0$. So, $b=1\pm \sqrt{2}$. Which, these provide $\lambda_\pm = \epsilon(b-1) = \pm \epsilon\sqrt{2}$, each corresponds to eigenvectors $\ket{\lambda_\pm} = \ket{1}+(1\pm\sqrt{2})\ket{2}$ (up to nonzero scaling).
        
        \hfil

        \item Under ordered basis $\alpha = \{\ket{1},\ket{2}\}$, $\hat{H}$ is given as:
        \begin{align}
            \mathcal{M}(\hat{H},\alpha)=\epsilon\begin{pmatrix}
                -1&1\\1&1
            \end{pmatrix}
        \end{align}
        If written in the ordered eigenvector basis $\beta=\{\ket{\lambda_+},\ket{\lambda_-}\}$, then $\hat{H}$ is represented as:
        \begin{align}
            \mathcal{M}(\hat{H},\beta)=\begin{pmatrix}
                \epsilon\sqrt{2} &0\\0&-\epsilon\sqrt{2}
            \end{pmatrix}
        \end{align}
    \end{enumerate}
\end{proof}

\newpage

\section*{10}
\begin{ques}\label{q10}
10 A three-level system (Griffiths 3.44)

\[
H = \begin{pmatrix} a & 0 & b \\ 0 & c & 0 \\ b & 0 & a \end{pmatrix}
\]

1. If the initial state is $\begin{pmatrix} 0 \\ 1 \\ 0\end{pmatrix}$, find $|S(t)\rangle$.

2. If the initial state is $\begin{pmatrix} 1 \\ 0 \\ 0\end{pmatrix}$, find $|S(t)\rangle$.
\end{ques}

\begin{proof}
    Given the hamiltonian as above, given the state $\ket{S(t)} = \begin{pmatrix}
        x(t)\\y(t)\\z(t)
    \end{pmatrix}$, the time-dependent Schr\"odinger equation states:
    \begin{align}
        \begin{pmatrix}
            x'(t)\\y'(t)\\z'(t)
        \end{pmatrix}&=\frac{\partial}{\partial t}\ket{S(t)}=\frac{1}{i\hbar}\hat{H}\ket{S(t)} = \begin{pmatrix}
            a/(i\hbar)&0&b/(i\hbar)\\0&c/(i\hbar)&0\\b/(i\hbar)&0&a/(i\hbar)
        \end{pmatrix}\begin{pmatrix}
            x(t)\\y(t)\\z(t)
        \end{pmatrix}
    \end{align}
    Which, compute the characeristic polynomial of the matrix $\frac{1}{i\hbar}\hat{H}$, one has:
    \begin{align}
        \det(\hat{H}/(i\hbar)-\lambda I) &= \left(\frac{c}{i\hbar}-\lambda\right)\left(\left(\frac{a}{i\hbar}-\lambda\right)^2-\left(\frac{b}{i\hbar}\right)^2\right)\\
        &= \left(\frac{c}{i\hbar}-\lambda\right)\left(\lambda^2 - \frac{2a}{i\hbar}\lambda + \frac{b^2-a^2}{\hbar^2}\right)
    \end{align}
    Which, the roots of the characteristic polynomials (or, the eigenvalues) includes $\lambda_1 = \frac{c}{i\hbar}$, and the other two are $\lambda_{2\pm} = \frac{2a/(i\hbar) \pm \sqrt{-4a^2/\hbar^2 + 4(a^2-b^2)/\hbar^2}}{2} = \frac{a\pm b}{i\hbar}$.

    \hfil

    Here, for simplicity assume $a,b,c\in\RR$ satisfy $(a+b),(a-b),c$ are distinct (so the three eigenvalues are distinct). Then, the eigenvectors corresponding to each eigenvalue are:
    \begin{align}
        \lambda_1:\quad (\hat{H}/(i\hbar) - \lambda_1 I)\begin{pmatrix}
            x\\y\\z
        \end{pmatrix} = 0 &\implies \frac{1}{i\hbar}\begin{pmatrix}
            a-c&0&b\\0&0&0\\b&0&a-c
        \end{pmatrix}\begin{pmatrix}
            x\\y\\z
        \end{pmatrix}=0 \implies \begin{pmatrix}
            1&0&0\\0&0&0\\0&0&1
        \end{pmatrix}\begin{pmatrix}
            x\\y\\z
        \end{pmatrix}=0\\
        &\implies x,z=0 \implies \begin{pmatrix}
            x\\y\\z
        \end{pmatrix} = y\begin{pmatrix}
            0\\1\\0
        \end{pmatrix}
    \end{align}

    \begin{align}
        \lambda_{2+}: \quad (\hat{H}/(i\hbar) - \lambda_{2+} I) \begin{pmatrix}
            x\\y\\z
        \end{pmatrix} = 0&\implies \frac{1}{i\hbar}\begin{pmatrix}
            -b&0&b\\0&c-(a+b)&0\\b&0&-b
        \end{pmatrix}\begin{pmatrix}
            x\\y\\z
        \end{pmatrix} = 0\implies \begin{pmatrix}
            1&0&-1\\0&1&0\\0&0&0
        \end{pmatrix}\begin{pmatrix}
            x\\y\\z
        \end{pmatrix}=0\\
        &\implies y=0,\ x=z \implies \begin{pmatrix}
            x\\y\\z
        \end{pmatrix} = z\begin{pmatrix}
            1\\0\\1
        \end{pmatrix}
    \end{align}

    \begin{align}
        \lambda_{2-}: \quad (\hat{H}/(i\hbar)-\lambda_{2-} I)\begin{pmatrix}
            x\\y\\z
        \end{pmatrix} = 0&=implies \frac{1}{i\hbar}\begin{pmatrix}
            b&0&b\\0&c-(a-b)&0\\b&0&b
        \end{pmatrix}\begin{pmatrix}
            x\\y\\z
        \end{pmatrix}=0\implies \begin{pmatrix}
            1&0&1\\0&1&0\\0&0&0
        \end{pmatrix}\begin{pmatrix}
            x\\y\\z
        \end{pmatrix}=0\\
        &\implies y=0,\ x=-z \implies \begin{pmatrix}
            x\\y\\z
        \end{pmatrix} = z\begin{pmatrix}
            -1\\0\\1
        \end{pmatrix}
    \end{align}
    Hence, the general solution to the given Schr\"odinger equation is spanned by the following three eigenvectors:
    \begin{align}
        \ket{\lambda_1}=e^{\lambda_1 t}\begin{pmatrix}
            0\\1\\0
        \end{pmatrix} = e^{-\frac{ict}{\hbar}}\begin{pmatrix}
            0\\1\\0
        \end{pmatrix},\quad \ket{\lambda_{2+}}=e^{\lambda_{2+}t}\begin{pmatrix}
            1\\0\\1
        \end{pmatrix} = e^{-\frac{i(a+b)t}{\hbar}}\begin{pmatrix}
            1\\0\\1
        \end{pmatrix},\quad \ket{\lambda_{2-}}=e^{\lambda_{2-}t}\begin{pmatrix}
            -1\\0\\1
        \end{pmatrix} = e^{-\frac{i(a-b)t}{\hbar}}\begin{pmatrix}
            -1\\0\\1
        \end{pmatrix}
    \end{align}
    Which, any state $\ket{S(t)} = a_1 \ket{\lambda_1}+a_{2+}\ket{\lambda_{2+}}+a_{2-}\ket{\lambda_{2-}}$ for some coefficients in $\CC$.

    \hfil

    \begin{enumerate}
        \item If the initial state is $\begin{pmatrix}
            0\\1\\0
        \end{pmatrix}$, plugin $t=0$, one yields $\ket{S(0)} = \begin{pmatrix}
            a_{2+}-a_{2-}\\a_1\\a_{2+}+a_{2-}
        \end{pmatrix}$, which enforces $a_1=1$, $a_{2+}-a_{2-}=0$ (or $a_{2+}=a_{2-}$), and $a_{2+}+a_{2-}=0$ (or $a_{2+}=-a-{2-}$). This enforces $a_{2+}=a_{2-}=0$. So, $\ket{S(t)} = \ket{\lambda_1}$, which is as follow:
        \begin{align}
            \ket{S(t)}=e^{-\frac{ict}{\hbar}}\begin{pmatrix}
                0\\1\\0
            \end{pmatrix}
        \end{align}

        \hfil

        \item If the initial state is $\begin{pmatrix}
            1\\0\\0
        \end{pmatrix}$, we again have $\ket{S(0)} = \begin{pmatrix}
            a_{2+}-a_{2-}\\a_1\\a_{2+}+a_{2-}
        \end{pmatrix}$, which enforces $a_1=0$, $a_{2+}+a_{2-}=0$ (or $a_{2-}=-a_{2+}$), and $a_{2+}-a_{2-}=1$ (which enforces $2a_{2+}=1$, or $a_{2+}=\frac{1}{2}$, and $a_{2-}=-\frac{1}{2}$). Hence, $\ket{S(t)} = \frac{1}{2}\ket{\lambda_{2+}} - \frac{1}{2}\ket{\lambda_-}$, which is as follow:
        \begin{align}
            \ket{S(t)} = \frac{1}{2}e^{-\frac{i(a+b)t}{\hbar}}\begin{pmatrix}
                1\\0\\1
            \end{pmatrix} - \frac{1}{2}e^{-\frac{i(a-b)t}{\hbar}}\begin{pmatrix}
                -1\\0\\1
            \end{pmatrix} = e^{-\frac{iat}{\hbar}}\begin{pmatrix}
                \frac{e^{ibt/\hbar}+e^{-ibt/\hbar}}{2}\\0\\-\frac{e^{ibt/\hbar}-e^{-ibt/\hbar}}{2} 
            \end{pmatrix}= e^{-iat/\hbar}\begin{pmatrix}
                    \cos(bt/\hbar)\\0\\-i\sin(bt/\hbar)
                \end{pmatrix}
        \end{align}
    \end{enumerate}
\end{proof}


\end{document}
